% Options for packages loaded elsewhere
\PassOptionsToPackage{unicode}{hyperref}
\PassOptionsToPackage{hyphens}{url}
\documentclass[
  a4paper,
]{ltjsarticle}
\usepackage{xcolor}
\usepackage[margin=25mm]{geometry}
\usepackage{amsmath,amssymb}
\setcounter{secnumdepth}{-\maxdimen} % remove section numbering
\usepackage{iftex}
\ifPDFTeX
  \usepackage[T1]{fontenc}
  \usepackage[utf8]{inputenc}
  \usepackage{textcomp} % provide euro and other symbols
\else % if luatex or xetex
  \usepackage{unicode-math} % this also loads fontspec
  \defaultfontfeatures{Scale=MatchLowercase}
  \defaultfontfeatures[\rmfamily]{Ligatures=TeX,Scale=1}
\fi
\usepackage{lmodern}
\ifPDFTeX\else
  % xetex/luatex font selection
\fi
% Use upquote if available, for straight quotes in verbatim environments
\IfFileExists{upquote.sty}{\usepackage{upquote}}{}
\IfFileExists{microtype.sty}{% use microtype if available
  \usepackage[]{microtype}
  \UseMicrotypeSet[protrusion]{basicmath} % disable protrusion for tt fonts
}{}
\makeatletter
\@ifundefined{KOMAClassName}{% if non-KOMA class
  \IfFileExists{parskip.sty}{%
    \usepackage{parskip}
  }{% else
    \setlength{\parindent}{0pt}
    \setlength{\parskip}{6pt plus 2pt minus 1pt}}
}{% if KOMA class
  \KOMAoptions{parskip=half}}
\makeatother
\usepackage{longtable,booktabs,array}
\newcounter{none} % for unnumbered tables
\usepackage{calc} % for calculating minipage widths
% Correct order of tables after \paragraph or \subparagraph
\usepackage{etoolbox}
\makeatletter
\patchcmd\longtable{\par}{\if@noskipsec\mbox{}\fi\par}{}{}
\makeatother
% Allow footnotes in longtable head/foot
\IfFileExists{footnotehyper.sty}{\usepackage{footnotehyper}}{\usepackage{footnote}}
\makesavenoteenv{longtable}
\setlength{\emergencystretch}{3em} % prevent overfull lines
\providecommand{\tightlist}{%
  \setlength{\itemsep}{0pt}\setlength{\parskip}{0pt}}
% Glyph map for lualatex (newunicodechar). Maps symbols that may lack font glyphs.
\usepackage{newunicodechar}
\usepackage{amssymb}
\newunicodechar{∎}{\ensuremath{\blacksquare}}
\newunicodechar{✓}{\ensuremath{\checkmark}}
\newunicodechar{✗}{\ensuremath{\times}}
\newunicodechar{⟹}{\ensuremath{\Longrightarrow}}
\newunicodechar{⟺}{\ensuremath{\Longleftrightarrow}}
\newunicodechar{↪}{\ensuremath{\hookrightarrow}}
\newunicodechar{⌊}{\ensuremath{\lfloor}}
\newunicodechar{⌋}{\ensuremath{\rfloor}}
\newunicodechar{≃}{\ensuremath{\simeq}}
\newunicodechar{≈}{\ensuremath{\approx}}
\newunicodechar{≤}{\ensuremath{\leq}}
\newunicodechar{≥}{\ensuremath{\geq}}
\newunicodechar{≠}{\ensuremath{\neq}}
\newunicodechar{→}{\ensuremath{\rightarrow}}
\newunicodechar{←}{\ensuremath{\leftarrow}}
\newunicodechar{↔}{\ensuremath{\leftrightarrow}}
\newunicodechar{⊥}{\ensuremath{\perp}}
\newunicodechar{√}{\ensuremath{\surd}}
\newunicodechar{∑}{\ensuremath{\sum}}
\newunicodechar{∏}{\ensuremath{\prod}}
\newunicodechar{∞}{\ensuremath{\infty}}
\newunicodechar{∈}{\ensuremath{\in}}
\newunicodechar{∀}{\ensuremath{\forall}}
\newunicodechar{∣}{\ensuremath{\mid}}
\newunicodechar{γ}{\ensuremath{\gamma}}
\newunicodechar{λ}{\ensuremath{\lambda}}
\newunicodechar{μ}{\ensuremath{\mu}}
\newunicodechar{ρ}{\ensuremath{\rho}}
\newunicodechar{σ}{\ensuremath{\sigma}}
\newunicodechar{φ}{\ensuremath{\varphi}}
\newunicodechar{ψ}{\ensuremath{\psi}}
\newunicodechar{θ}{\ensuremath{\theta}}
\newunicodechar{Δ}{\ensuremath{\Delta}}
\newunicodechar{Π}{\ensuremath{\Pi}}
\newunicodechar{Λ}{\ensuremath{\Lambda}}
\newunicodechar{τ}{\ensuremath{\tau}}
\newunicodechar{ε}{\ensuremath{\varepsilon}}
\newunicodechar{ω}{\ensuremath{\omega}}
\newunicodechar{π}{\ensuremath{\pi}}
\newunicodechar{ν}{\ensuremath{\nu}}
\newunicodechar{±}{\ensuremath{\pm}}
\newunicodechar{×}{\ensuremath{\times}}
\newunicodechar{·}{\ensuremath{\cdot}}
\usepackage{bookmark}
\IfFileExists{xurl.sty}{\usepackage{xurl}}{} % add URL line breaks if available
\urlstyle{same}
\hypersetup{
  pdftitle={Anonymity and the Minimal Unit of Information --- The minimal zero-closure unit is 2 waves; how many waves is the minimal unit of information that survives after quotienting by S\_N? (Hypothesis / Exploratory Paper)},
  pdfauthor={Noriaki Kihara (WF System Co., Ltd.), ORCID 0009-0004-6753-4020},
  hidelinks,
  pdfcreator={LaTeX via pandoc}}

\title{Anonymity and the Minimal Unit of Information --- The minimal
zero-closure unit is 2 waves; how many waves is the minimal unit of
information that survives after quotienting by \(S_N\)? (Hypothesis /
Exploratory Paper)}
\author{Noriaki Kihara (WF System Co., Ltd.), ORCID 0009-0004-6753-4020}
\date{2026-09-12}

\begin{document}
\maketitle

\textbf{Series:} Foundations of Self-Consistent Relational-Wave Closed
Systems (Paper 9 of 10 / Exploratory Paper)\\
\textbf{Author:} Noriaki Kihara (WF System Co., Ltd.)　\textbf{Date:}
2026-09-12\\
\textbf{Version DOI:} 10.5281/zenodo.22729128\\
\textbf{Concept DOI:} 10.5281/zenodo.22729127\\
\textbf{ORCID:} 0009-0004-6753-4020\\
\textbf{Zenodo:} https://zenodo.org/records/22729128\\

\begin{center}\rule{0.5\linewidth}{0.5pt}\end{center}

\subsection{Abstract}\label{abstract}

This paper is a \textbf{hypothesis / exploratory paper} (no new runs; it
organizes algebraic facts + existing empirical data + literature
cross-checks, with claim classification made explicit). We examine, from
the standpoint of anonymity, ``what can carry information'' in a
relational-wave closed system. (1) The edge index \(m\) is a gauge of
the vertex permutation \(S_N\), and ``the type of wave number \(m\)'' is
not a physical question (mathematical fact). (2) The final state is
exactly equimodular (relative spread \(N=12{:}1.17\times10^{-13}\),
\(N=40{:}3.54\times10^{-14}\)), and under 12-digit rounding there is
only one kind of \(|z|\) --- indistinguishable even by value, and the
indistinguishability of identical particles appears as a consequence of
the dynamics (numerical verification). (3) The nontrivial solution of
closure \(\sum z^2=0\) is minimally 2 waves (\(z_2=\pm i z_1\)); a
single wave admits only \(z=0\) and cannot carry information
(mathematical fact = \textbf{minimal closure unit is 2 waves}). However,
\(\pm\) flips under pair exchange (\(S_2\)), so on its own it is not an
\(S_N\)-invariant bit. The primitive block size is floor-dependent,
being \(p=\operatorname{spf}(L)\) of Paper 4 (\(L=N/\gcd(N,4)\); only
\(8\mid N\) reaches the algebraic lower bound 2). Hence \textbf{the
carrier of information is not the individual wave but the minimal
zero-closure block} (Lam--Leung); but ``2 waves = minimal closure unit''
does not mean ``2 waves = anonymous information unit.'' Analyzing
\(\mathcal C_k\simeq Q_{k-2}/S_k\) (projective quadric hypersurface
\(\dim_{\mathbb C}=k-2\)), \textbf{closure alone does not generate
discrete anonymous species} (\(k=2\) collapses to 1 orbit under \(S_2\);
\(k\ge3\) is a continuous moduli). The discrete invariant is not \(\pm\)
but the block degree \(p=\operatorname{spf}(L)\) (Paper 4; the closure
axiom + the discrete symmetry of the floor quantizes it), and the
self-consistent dynamics that quantizes the continuous closure moduli
into discrete types is the next question (§5). (4) A structure that
records anything beyond \(a,b\) on a wave (the band/hair register of
prior work) is a \textbf{scaffold} not derivable from the closure axiom
or from anonymity, and the types read there (boson/fermion) are not
cited as a physical claim (downgrade of the verdict). Classification: we
distinguish mathematical fact / numerical verification / hypothesis.

\begin{center}\rule{0.5\linewidth}{0.5pt}\end{center}

\subsection{1. Problem}\label{problem}

The waves of this system are anonymous (they carry no external index).
We examine, using algebra and existing empirical data, what information
can reside in an anonymous object and whether distinctions such as
particle species can be made at all. This paper does not assert
conclusions; it classifies and fixes what is an established fact and
what is a hypothesis.

\subsection{\texorpdfstring{2. The edge index is an \(S_N\) gauge
(mathematical
fact)}{2. The edge index is an S\_N gauge (mathematical fact)}}\label{the-edge-index-is-an-s_n-gauge-mathematical-fact}

In the generator \(K_{ef}=A_{ef}\sin(\varphi_f-\varphi_e)\), \(A\) is
the edge-graph adjacency of \(K_N\), and the dynamics is equivariant
under the vertex permutation \(S_N\). The numbering of edges is a
product of the choice of vertex labels; permuting them moves to a
different wave. Therefore ``the type of wave number \(m\)'' is not a
physical question, and \textbf{what is readable are the
\(S_N\)-invariants} (totals, multisets, the count ``how many waves of
type \(X\)''). What is not readable is ``which wave'' (identification /
tracking).

This is the \textbf{same principle} as the relational readout of Paper 7
(the absolute normal cannot be recovered from the state alone; what is
readable is the relative phase between faces) --- physical readout
resides not in absolute labels (which wave / which orientation) but in
relational quantities invariant under transformation. Adding Paper 8:
(i) \textbf{kinematic anonymity} = the wave index is an \(S_N\) gauge
from the outset, (ii) \textbf{dynamical anonymization} = at termination
even the \(|z_e|\) become equal so identification by value also
vanishes, (iii) \textbf{what remains} = the relational invariants of
phase relations, relative arrangement of faces, and the
type/multiplicity of closure blocks. This is the relationalism of the
whole series.

\subsection{3. The final state is indistinguishable even by value
(numerical
verification)}\label{the-final-state-is-indistinguishable-even-by-value-numerical-verification}

The final state is exactly equimodular: relative spread
\(1.17\times10^{-13}\) for \(N=12\) and \(3.54\times10^{-14}\) for
\(N=40\). Under 12-digit rounding there is only one kind of \(|z|\)
(\(66\to1\), \(780\to1\)). At step0 they are all distinct (\(66/66\),
\(780/780\)), so the anonymization is \textbf{driven by the dynamics}.
That is, the indistinguishability of identical particles appears not as
an initial condition but as a consequence of the dynamics (a different
expression of the same fact as the equal amplitude of Paper 5 and
\(a_\infty=\sqrt{H/M}\) of Paper 8, and consistent with the anonymity
H-theorem \(S\to1\) of Experiment 7).

\subsection{4. The minimal unit of information is 2 waves (mathematical
fact +
hypothesis)}\label{the-minimal-unit-of-information-is-2-waves-mathematical-fact-hypothesis}

Nontrivial solutions of closure \(\sum z^2=0\):

{\def\LTcaptype{none} % do not increment counter
\begin{longtable}[]{@{}
  >{\raggedright\arraybackslash}p{(\linewidth - 4\tabcolsep) * \real{0.3333}}
  >{\raggedright\arraybackslash}p{(\linewidth - 4\tabcolsep) * \real{0.3333}}
  >{\raggedright\arraybackslash}p{(\linewidth - 4\tabcolsep) * \real{0.3333}}@{}}
\toprule\noalign{}
\begin{minipage}[b]{\linewidth}\raggedright
Wave count
\end{minipage} & \begin{minipage}[b]{\linewidth}\raggedright
Solution
\end{minipage} & \begin{minipage}[b]{\linewidth}\raggedright
Degrees of freedom
\end{minipage} \\
\midrule\noalign{}
\endhead
\bottomrule\noalign{}
\endlastfoot
1 & \(z^2=0\Rightarrow z=0\) only & None = cannot carry information \\
2 & \(z_2=\pm i z_1\) (minimal nontrivial closure unit) & all of \(z_1\)
(scale, phase). The sign \(\pm\) flips under pair exchange (see
below) \\
3 & cube-root type \(z_k=z\,\omega^k\), etc. & wider solution variety \\
\end{longtable}
}

\textbf{The minimal closure unit is 2 waves (mathematical fact)}: a
nonzero closure block requires at least 2 waves; a single wave admits
only \(z=0\) and cannot carry information.

\textbf{However, \(\pm\) on its own is not an \(S_N\)-invariant 2-valued
quantity (mathematical fact)}: pair-exchanging (\(S_2\subset S_N\)) the
pair \(z_2=+iz_1\) gives \(z'_1=z_2,\ z'_2=z_1\) so that \(z'_2=-iz'_1\)
--- \textbf{the very \(S_2\) action that renders the two waves of the
pair anonymous flips the sign}. Hence \(+i\leftrightarrow-i\) lie on the
same \(S_N\) orbit, and to physicalize a two-valued quantity such as
right-/left-handed one needs additional relational structure that can
compare the signs (propagation direction, orientation, a third relation;
the helicity discussion of Paper 4). That is, \textbf{the existence of a
closure unit \(\neq\) the existence of an anonymous information bit}.

\textbf{The primitive block size is floor-dependent (Paper 4)}: by the
general rule \(L=N/\gcd(N,4)\), \(p=\operatorname{spf}(L)\) of Paper 4,
the primitive block takes \(p=2,3,5,7,\ldots\)
(\(N=10\Rightarrow L=5\Rightarrow p=5\); odd prime \(N\Rightarrow p=N\);
Lam--Leung). Therefore

\[k_{\min}=2\ (\text{absolute algebraic lower bound})\qquad\text{vs}\qquad k_{\rm primitive}=\operatorname{spf}(L)\ (\text{the minimal unit realized on that floor}),\]

the ``minimal \textbf{possible} unit'' and the ``minimal unit actually
allowed'' differ, and only \(8\mid N\) reaches the lower bound 2.

\textbf{Consequence (hypothesis) and the next question}: the carrier of
information is not the individual wave but the minimal closure block (a
single wave cannot close and cannot carry meaning). However, by the
above, ``2 waves = minimal closure unit'' does not mean ``2 waves =
anonymous information unit.'' What should really be asked is the minimal
\(k\) such that discretely distinct types still remain after quotienting
the closure solutions by symmetry:

\[\mathcal C_k=\Big\{(z_1,\ldots,z_k):\ \textstyle\sum_{j=1}^k z_j^2=0\Big\}\Big/(\text{global phase, scale, }S_k)\]

what is the minimal \(k\) such that discretely distinct types remain (=
\textbf{minimal anonymous information unit})? \(k_{\min}=2\) is the
minimal closure unit. The minimal anonymous information unit is analyzed
in §5 (\(\mathcal C_k\simeq Q_{k-2}/S_k\)). Candidate particle species
(block size, multiplicity, way of assembly) can be defined as
\(S_N\)-invariant multisets and do not break anonymity, but \(\pm\) on
its own is not an invariant bit.

\subsection{\texorpdfstring{5. The closure moduli are continuous ---
discrete anonymous species require additional self-consistent structure
(\(\mathcal C_k\simeq Q_{k-2}/S_k\))}{5. The closure moduli are continuous --- discrete anonymous species require additional self-consistent structure (\textbackslash mathcal C\_k\textbackslash simeq Q\_\{k-2\}/S\_k)}}\label{the-closure-moduli-are-continuous-discrete-anonymous-species-require-additional-self-consistent-structure-mathcal-c_ksimeq-q_k-2s_k}

The \(\mathcal C_k\) of §4 can be analyzed without further computation.
Quotienting phase and scale together by \(z\sim cz\)
(\(c\in\mathbb C^\times\)), \(\sum_{j=1}^k z_j^2=0\) becomes a
\textbf{projective quadric hypersurface} on \(\mathbb{CP}^{k-1}\),

\[Q_{k-2}=\Big\{[z_1:\cdots:z_k]\in\mathbb{CP}^{k-1}:\ \textstyle\sum_j z_j^2=0\Big\},\qquad \dim_{\mathbb C}Q_{k-2}=k-2,\]

and \(\mathcal C_k\simeq Q_{k-2}/S_k\).

\begin{itemize}
\tightlist
\item
  \textbf{\(k=2\)}: the projective solutions of \(z_1^2+z_2^2=0\) are
  the 2 points \([1:i],[1:-i]\) (\(\dim_{\mathbb C}Q_0=0\)), but \(S_2\)
  exchanges the two, so \(Q_0/S_2\) is \textbf{1 orbit} --- \textbf{for
  2 waves there is closure but no anonymous 2-valued type} (an
  alternative proof of the \(\pm\) flip of §4).
\item
  \textbf{\(k\ge3\)}: \(\dim_{\mathbb C}Q_{k-2}=k-2\ge1\). Since \(S_k\)
  is a finite group, taking the quotient does not remove the continuous
  dimension near a generic point, and
  \(\dim_{\mathbb C}(Q_{k-2}/S_k)=k-2\) --- the solution space is still
  a \textbf{continuous moduli}.
\end{itemize}

Therefore \textbf{the closure condition \(\sum z^2=0\) alone does not,
after quotienting by \(S_k\), produce a finite number of discretely
distinct types (particle species)} (\(k=2\) collapses to 1 species under
\(S_2\); \(k\ge3\) retains continuous degrees of freedom). Hence the
question shifts from ``how many waves can carry information'' to

\[\boxed{\text{which self-consistent dynamics quantizes the continuous closure moduli into discrete types}}\]

\textbf{Connection with Paper 4 (the carrier of quantization)}: in Paper
4, the \textbf{discrete block size}
\(p=\operatorname{spf}(N/\gcd(N,4))\) emerged from the cyclic structure
of the high-symmetry floor. Although the closure equation alone is a
continuous moduli, \textbf{imposing the discrete symmetry of the
high-symmetry floor quantizes the block size}:

\[\text{closure axiom}\ +\ \text{discrete symmetry of the floor}\ \Longrightarrow\ \text{quantization of block size }p=2,3,5,7,\ldots.\]

At the present stage the invariant that is reliably discrete is not
\(\pm\) but the \textbf{block degree \(p\) itself}, and whether a
discrete invariant corresponding to a further particle species arises
within it is the next problem. Rather than adding, from outside the
closure, the register of Paper A (128 declared constants), this paper
points in the direction that discrete types must be generated
\textbf{from within} by closure + floor + dynamical selection rules
(orientation, winding, topology, etc.).

\subsection{\texorpdfstring{6. Recording anything beyond \(a,b\) on a
wave is an anonymity violation (downgrade of the
verdict)}{6. Recording anything beyond a,b on a wave is an anonymity violation (downgrade of the verdict)}}\label{recording-anything-beyond-ab-on-a-wave-is-an-anonymity-violation-downgrade-of-the-verdict}

\begin{itemize}
\tightlist
\item
  The waves of this system are a single complex number
  (\(z\in\mathbb C^M\), only \(a+ib\)) = the minimal description.
\item
  Prior work (Paper A) assigns each wave \(16\times8=128\) complex
  numbers (band \(k\) / hair \(\eta\) registers). \(N_n=16,N_\eta=8\)
  are \textbf{declared constants}, derived neither from \(N\) nor from
  the closure axiom \(\sum z^2=0\). A wave becomes distinguishable by
  internal structure = anonymity is broken from outside.
\item
  Therefore the types read within that register (the boson/fermion
  distinction, ``the clock is ticked by fermions,'' ``charge = closure
  defect'') are \textbf{properties of the scaffold, not physics derived
  from the closure axiom}. This series does not cite these as physical
  claims (downgrade of the verdict). To treat the types as physics, one
  must derive the band-register-equivalent degrees of freedom from the
  closure axiom and anonymity (currently not achieved).
\end{itemize}

\subsection{7. Claim classification /
open}\label{claim-classification-open}

\begin{itemize}
\tightlist
\item
  \textbf{Mathematical fact}: the \(S_N\) gauge (§2), the minimal
  nontrivial closure unit of \(\sum z^2=0\) = 2 waves \(z_2=\pm i z_1\)
  (§4; but \(\pm\) flips under pair exchange \(S_2\) and on its own is
  not an \(S_N\)-invariant bit), the primitive block size
  \(=\operatorname{spf}(L)\) (\(L=N/\gcd(N,4)\); only \(8\mid N\)
  reaches the algebraic lower bound 2, §4 · Paper 4), the Lam--Leung
  decomposition.
\item
  \textbf{Numerical verification}: the final state is exactly
  equimodular (§3).
\item
  \textbf{Hypothesis}: the carrier of information = the minimal closure
  block (not the individual wave), particle species = \(S_N\)-invariant
  multisets (block size, multiplicity, way of assembly) (§4).
\item
  Verified (Reproduction B): the test of the block decomposition of the
  final state was carried out for all \(N=3\)--\(40\) (existing
  10000-step terminal states). The terminal state is exactly equimodular
  (deviation \(\sim10^{-13}\)) and conserves \(\Sigma z^2\approx0\), but
  \textbf{exact \(\pm i\) pairs are zero} (the make\_parent floor as
  well; the minimal \(\varepsilon\) only decreases asymptotically to
  \(\sim2\times10^{-6}\) at large \(N\)). That is, minimal-block
  (\(\pm i\) pair) decomposition is \textbf{specific to the
  cyclic-Fourier / integer-K high-symmetry floor (\(N=8k\), Paper 4)}
  and does not stand at the terminus of the data series. Hence
  ``terminal block decomposition = catalog of species'' does not hold as
  a property of the terminal state, and block structure (the type
  corresponding to species) resides on the high-symmetry-floor side
  (Paper 4) --- consistent with the thrust of this paper (the carrier of
  information is the block; the type compatible with anonymity resides
  in the closure structure of the floor).
\item
  \textbf{Analytic answer (§5 · mathematical fact)}: closure
  \(\sum z^2=0\) alone does not generate discrete anonymous species ---
  with \(\mathcal C_k\simeq Q_{k-2}/S_k\) (projective quadric
  hypersurface \(\dim_{\mathbb C}Q_{k-2}=k-2\)), \(k=2\) collapses to 1
  orbit under \(S_2\), and \(k\ge3\) is a continuous moduli
  (\(\dim_{\mathbb C}=k-2\ge1\), not removed by the finite group
  \(S_k\)). What is discrete is not \(\pm\) but the \textbf{block degree
  \(p=\operatorname{spf}(N/\gcd(N,4))\)} (Paper 4), where the closure
  axiom + the discrete symmetry of the floor quantizes the block size.
\item
  Open / next research question: what is the self-consistent dynamics
  that quantizes the continuous closure moduli into discrete types
  (dynamical selection rules, orientation, winding, topology, etc.
  should generate them from within)? Does a discrete invariant
  corresponding to a further particle species arise within the block
  degree \(p\)? Reproduction B showed that exact \(\pm i\) pairs
  (\(p=2\) blocks) are specific to the high-symmetry floor \(8\mid N\)
  and vanish at the generic terminus (the type corresponding to species
  resides not in the individual wave but on the block side, compatible
  with anonymity).
\end{itemize}

\subsection{Related work (structural agreement, not the source of
derivation)}\label{related-work-structural-agreement-not-the-source-of-derivation}

\begin{itemize}
\tightlist
\item
  Vanishing sums (Lam--Leung): prime-size decomposition of zero-closure
  blocks.
\item
  Indistinguishability of identical particles (quantum statistics): in
  this system it appears as a consequence of the dynamics.
\end{itemize}

\subsection{References}\label{references}

\begin{enumerate}
\def\labelenumi{\arabic{enumi}.}
\tightlist
\item
  Noriaki Kihara, ``The mechanism of inflationary rapid expansion in
  self-consistent relational-wave closed systems,'' Concept DOI
  10.5281/zenodo.22112008.
\item
  Noriaki Kihara, ``Conditions under which a seed gives rise to
  particles and the lower bound of resolution (Paper A),'' Concept DOI
  10.5281/zenodo.21874481. (Source of the band/hair register referenced
  in §5.)
\item
  T. Y. Lam and K. H. Leung, ``On vanishing sums of roots of unity,'' J.
  Algebra 224 (2000) 91.
\end{enumerate}

\subsection{Reproduction}\label{reproduction}

This paper is a hypothesis / exploratory paper with no new runs. The
numerical verification (final-state equimodularity) relies on the
existing data \texttt{N3\_N40\_long10000\_20260905/} (10000-step sweep,
\texttt{check\_final\_states\_10000\_v1.py}, etc.). The \(N\)-dependence
test of the terminal block decomposition (Reproduction B) is in
\texttt{位相ロック全N確認\_相対平衡\_20260912/終端状態ブロック分解\_N依存\_20260912/}
(\texttt{analyze\_terminal\_block\_decomposition\_20260912.py},
\texttt{results/terminal\_block\_summary.csv}, README / REPORT / SHA /
run\_all; read-only). The closure factorization is shared with the
package of Paper 4. The algebra (the minimal solution of \(\sum z^2=0\),
Lam--Leung) is self-contained in the main text. The primary record of
the study is in the discovery ledger
\texttt{発見台帳\_干渉保存力学\_20260905起点.md} (Studies 25, 26, 27).

\end{document}
